\chapter*{术语与状态标记}
\addcontentsline{toc}{chapter}{术语与状态标记}

\begin{longtable}{>{\raggedright\arraybackslash}p{4.0cm}>{\raggedright\arraybackslash}p{9.6cm}}
\caption{关键术语与证据状态}\label{tab:abbreviations}\\
\toprule
\textbf{术语} & \textbf{本文含义}\\
\midrule
主体（Subject） & 跨模型和运行进程保持逻辑身份的主体。\\
认知权威（Cognitive Authority） & 持久认知对象、修订和生命周期的权威数据及其写入权限。\\
运行时（Runtime） & 会话、活动认知和任务上下文等短中期状态。\\
检索服务视图（Serving Projection） & 可从权威数据重建的索引、向量、拓扑或检索视图。\\
证据（Evidence） & 指向观察事件、来源区域或派生表示的可追溯引用。\\
访问控制（ACL） & 访问控制列表或访问策略。\\
比较并交换（CAS） & 依据预期修订或版本进行并发写入校验。\\
EPA & 研究资料包中的语义线索感知候选方法。\\
个性化 PageRank（PPR） & 候选拓扑扩散基线。\\
倒数排序融合（RRF） & 按排名合并多个证据通道的方法。\\
\midrule
\StatusTag{必需}{必要能力或语义约束。}\\
\StatusTag{目标}{当前研究确定的目标设计。}\\
\StatusTag{已实现}{当前代码中可定位的实现路径。}\\
\StatusTag{已核验}{有日期、平台和范围记录的执行结果。}\\
\StatusTag{实验观察}{研究资料包记载的局部实验结果。}\\
\StatusTag{候选}{需与基线和替代方案比较的方案。}\\
\StatusTag{待定}{仍需收敛的语义、归属或运行合同。}\\
\bottomrule
\end{longtable}

状态标签描述对应能力或主张的研究和实现状态。\clearpage
